% ASSERT_CONVENTION: natural_units=natural, metric_signature=mostly_minus, fourier_convention=physics, coupling_convention=alpha_s, generator_normalization=T(fund)=1/2, covariant_derivative_sign=D_mu=partial_mu-igA
%
% Exercise Set 3: Seiberg Duality in Practice
% Part of: The Dualised Standard Model -- Part III Exercises
%
% Conventions:
%   R_fermion = R_superfield - 1 for chiral multiplet fermions
%   R_gaugino = +1 (directly)
%   B[Q] = +1, B[Q-tilde] = -1
%   B[q] = N_c/(N_f - N_c), B[q-tilde] = -N_c/(N_f - N_c)
%   A(fund) = 1, A(antifund) = -1
%   T(fund) = T(antifund) = 1/2
%
\documentclass[11pt,a4paper]{article}
% ASSERT_CONVENTION: natural_units=natural, metric_signature=mostly_minus, fourier_convention=physics, coupling_convention=alpha_s, generator_normalization=T(fund)=1/2, covariant_derivative_sign=D_mu=partial_mu-igA
%
% CONVENTIONS (Seiberg hep-th/9411149)
% Metric: (+,-,-,-) (mostly minus)
% T(fund) = 1/2 per Weyl fermion
% b_0 = (11N_c - 2N_f)/3 for N_f Dirac flavours
% D_mu = partial_mu - ig A_mu^a T^a
% A(fund) = 1 (cubic anomaly coefficient per fundamental Weyl fermion)
% alpha_s = g_s^2 / (4 pi)
% R[Q] = (N_f - N_c)/N_f; R[W] = 2
% N_f counts Dirac flavour pairs (each = 2 Weyl fermions)
%
% This file is \input{} by all lecture files.

%% ---------- Document class ----------
% (loaded by the wrapper; preamble only provides packages and macros)

%% ---------- Standard packages ----------
\usepackage{amsmath,amssymb,amsthm,mathtools}
\usepackage{hyperref}
\usepackage[capitalise,noabbrev]{cleveref}
\usepackage{enumitem}
\usepackage{booktabs}
\usepackage{array}
\usepackage{xcolor}
\usepackage{tikz}
\usetikzlibrary{arrows.meta,decorations.markings,positioning,calc}
\usepackage{fancybox}
\usepackage{tcolorbox}
\tcbuselibrary{skins,breakable}
\usepackage{slashed}              % Feynman slash notation
\usepackage{cite}                 % compressed citation lists

%% ---------- Hyperref configuration ----------
\hypersetup{
  colorlinks=true,
  linkcolor=blue!70!black,
  citecolor=green!50!black,
  urlcolor=blue!80!black,
  pdfauthor={Part III Lecture Notes},
  pdftitle={The Dualised Standard Model}
}

%% ---------- Theorem environments ----------
\theoremstyle{plain}
\newtheorem{theorem}{Theorem}[section]
\newtheorem{lemma}[theorem]{Lemma}
\newtheorem{proposition}[theorem]{Proposition}
\newtheorem{corollary}[theorem]{Corollary}

\theoremstyle{definition}
\newtheorem{definition}[theorem]{Definition}
\newtheorem{example}[theorem]{Example}
\newtheorem{exercise}{Exercise}[section]

\theoremstyle{remark}
\newtheorem{remark}[theorem]{Remark}

%% ---------- Physics macros: gauge groups and representations ----------
\newcommand{\Nc}{N_{\!c}}
\newcommand{\Nf}{N_{\!f}}
\newcommand{\SU}[1]{\mathrm{SU}(#1)}
\newcommand{\UU}[1]{\mathrm{U}(#1)}

% Representations
\newcommand{\fund}{\square}
\newcommand{\antifund}{\overline{\square}}
\newcommand{\adj}{\mathrm{adj}}
\newcommand{\antisym}{\wedge^{\!2}}        % antisymmetric rank-2
\newcommand{\sym}{\mathrm{S}_2}           % symmetric rank-2

%% ---------- Physics macros: group theory constants ----------
% Dynkin index: Tr(T^a T^b) = T(R) delta^{ab}
\newcommand{\Tfund}{\tfrac{1}{2}}          % T(fund) = 1/2 per Weyl fermion
\newcommand{\Tadj}{\Nc}                    % T(adj) = N_c
\newcommand{\Cfund}{\frac{\Nc^2 - 1}{2\Nc}}  % C_2(fund)
\newcommand{\Cadj}{\Nc}                    % C_2(adj)

% Cubic anomaly coefficient: A(R) = Tr_R[T^a {T^b, T^c}] with A(fund) = 1
\newcommand{\Afund}{1}

%% ---------- Physics macros: beta function ----------
\newcommand{\bzero}{b_0}
\newcommand{\bone}{b_1}
\newcommand{\alphas}{\alpha_s}
\newcommand{\gs}{g_s}

%% ---------- Physics macros: SUSY / R-charge ----------
\newcommand{\Rcharge}[1]{R[#1]}
\newcommand{\Wsup}{W}                      % superpotential

%% ---------- Physics macros: covariant derivative and field strength ----------
\newcommand{\Dmu}{D_\mu}
\newcommand{\Fmn}{F_{\mu\nu}}
\newcommand{\Ftilde}{\widetilde{F}_{\mu\nu}}

%% ---------- Physics macros: common abbreviations ----------
\newcommand{\ie}{\textit{i.e.}}
\newcommand{\eg}{\textit{e.g.}}
\newcommand{\cf}{\textit{cf.}}
\newcommand{\IR}{\ensuremath{\mathrm{IR}}}
\newcommand{\UV}{\ensuremath{\mathrm{UV}}}
\newcommand{\AF}{\ensuremath{\mathrm{AF}}}
\newcommand{\BZ}{\ensuremath{\mathrm{BZ}}}
\newcommand{\NSVZ}{\ensuremath{\mathrm{NSVZ}}}
\newcommand{\SQCD}{\ensuremath{\mathrm{SQCD}}}
\newcommand{\QCD}{\ensuremath{\mathrm{QCD}}}
\newcommand{\SM}{\ensuremath{\mathrm{SM}}}
\newcommand{\Tr}{\mathrm{Tr}}

%% ---------- Convention box macro ----------
\newtcolorbox{conventionbox}[1][]{%
  enhanced,
  colback=blue!5!white,
  colframe=blue!60!black,
  fonttitle=\bfseries,
  title={Conventions},
  #1
}

%% ---------- Comparison table style ----------
\newtcolorbox{comparisontable}[1][]{%
  enhanced,
  colback=orange!5!white,
  colframe=orange!70!black,
  fonttitle=\bfseries,
  title={Comparison},
  #1
}

%% ---------- Warning / important box ----------
\newtcolorbox{warningbox}[1][]{%
  enhanced,
  colback=red!5!white,
  colframe=red!70!black,
  fonttitle=\bfseries,
  title={Important},
  #1
}

%% ---------- Preview / forward-reference box ----------
\newtcolorbox{previewbox}[1][]{%
  enhanced,
  colback=green!5!white,
  colframe=green!50!black,
  fonttitle=\bfseries,
  title={Preview},
  #1
}

%% ---------- Page layout ----------
\usepackage[margin=2.5cm]{geometry}
\setlength{\parskip}{0.4em}
\setlength{\parindent}{0em}

%% ---------- Section formatting ----------
\usepackage{titlesec}
\titleformat{\section}{\Large\bfseries}{\thesection}{1em}{}
\titleformat{\subsection}{\large\bfseries}{\thesubsection}{1em}{}

%% ---------- End of preamble ----------


\title{\textbf{Exercise Set 3: Seiberg Duality in Practice}}
\author{The Dualised Standard Model --- Part III Exercises}
\date{Lent Term 2027}

\begin{document}
\maketitle

\begin{conventionbox}[title={Conventions for these exercises}]
\renewcommand{\arraystretch}{1.3}
\begin{tabular}{@{}l l@{}}
\toprule
\textbf{Quantity} & \textbf{Convention} \\
\midrule
Dynkin index & $T(\fund) = T(\antifund) = \Tfund$ \\
Cubic anomaly & $A(\fund) = 1$, \;$A(\antifund) = -1$ \\
$R$-charge (superfield) & $\Rcharge{Q} = (\Nf - \Nc)/\Nf$; \;$\Rcharge{\Wsup} = 2$ \\
$R$-charge (fermion) & $R_{\mathrm{ferm}} = R_{\mathrm{sf}} - 1$ (chiral);
  \;$R_{\lambda} = +1$ (gaugino) \\
Baryon charges & $B[Q] = +1$, \;$B[\widetilde{Q}] = -1$;\quad
  $B[q] = \Nc/(\Nf - \Nc)$ \\
$\Nf$ counting & Dirac pairs $(Q_i, \widetilde{Q}_i)$ \\
Magnetic theory & $\SU{\Nf - \Nc}$ with $(q, \widetilde{q}, M)$;
  \quad $\Wsup_{\mathrm{mag}} = \frac{1}{\mu}\, M\, q\, \widetilde{q}$ \\
\bottomrule
\end{tabular}

\medskip

\noindent References: Lecture~3 (\S4, anomaly matching); Lecture~4 (\S5, APS deformation).

\medskip

\noindent Solutions can be cross-checked against
\texttt{scripts/verify\_anomaly\_matching.py}.
\end{conventionbox}

\bigskip


%% ========================================================================
%%  PROBLEM 1: Anomaly Matching for SU(3) with N_f = 6
%% ========================================================================
\section*{Problem 1: Anomaly Matching for $\SU{3}$ with $\Nf = 6$
  \hfill $\star\star$ \normalsize(30 min)}
\addcontentsline{toc}{section}{Problem 1: Anomaly matching for SU(3),
  $\Nf = 6$}
\label{prob:anomaly-su3-6}

\noindent\textbf{Background.}
In Lecture~3, we verified all six independent 't~Hooft anomaly
coefficients as polynomial identities in $\Nc$ and $\Nf$.  In this
problem, you will carry out the explicit numerical computation for the
benchmark case $\Nc = 3$, $\Nf = 6$.

\medskip

\noindent\textbf{Setup.}  The electric theory is $\SU{3}$ with
$6 \times (Q, \widetilde{Q})$ and $\Wsup = 0$.  The magnetic dual
is $\SU{6 - 3} = \SU{3}$ with $6 \times (q, \widetilde{q})$,
gauge-singlet mesons $M^i{}_j$ ($6 \times 6$ matrix), and
$\Wsup_{\mathrm{mag}} = (1/\mu)\, M\, q\, \widetilde{q}$.

\medskip

\noindent\textbf{Questions.}
\begin{enumerate}[label=(\alph*)]
  \item \textbf{Quantum numbers.}  Write down the complete quantum number
    tables for both theories, including both superfield and fermion
    $R$-charges.  Use the values:
    \begin{align*}
      R_{\mathrm{sf}}[Q] &= (6-3)/6 = 1/2, &
      R_{\mathrm{ferm}}[Q] &= -1/2, \\
      R_{\mathrm{sf}}[q] &= 3/6 = 1/2, &
      R_{\mathrm{ferm}}[q] &= -1/2, \\
      R_{\mathrm{sf}}[M] &= 2(6-3)/6 = 1, &
      R_{\mathrm{ferm}}[M] &= 0, \\
      B[q] &= 3/(6-3) = 1. &
    \end{align*}

  \item \textbf{Compute all six anomaly coefficients} on both sides:
    \begin{enumerate}[label=(\roman*)]
      \item $\SU{6}_L^3$
      \item $\SU{6}_R^3$
      \item $\SU{6}_L^2\,\UU{1}_B$
      \item $\SU{6}_L^2\,\UU{1}_R$
      \item $\UU{1}_B^2\,\UU{1}_R$
      \item $\Tr[R_{\mathrm{ferm}}]$ (gravitational anomaly)
    \end{enumerate}

  \item Verify that all six coefficients match.
\end{enumerate}

\medskip

\noindent\textit{Hint:} For the mixed anomaly $\SU{\Nf}_L^2\,\UU{1}_X$,
the formula is $\mathcal{A} = \sum_f (\text{gauge mult}) \times
T(R^{(f)}_L) \times X_f$, summed over all fermions charged under
$\SU{\Nf}_L$.  Remember that only fermion (not superfield) $R$-charges
appear in anomaly triangles.

\bigskip
\hrule
\bigskip

\begin{proof}[\textbf{Solution to Problem 1}]

\textbf{(a) Quantum number tables.}

\textbf{Electric theory} ($\SU{3}$ gauge, $\Nf = 6$):

\renewcommand{\arraystretch}{1.3}
\begin{center}
\begin{tabular}{@{}l c c c c c c@{}}
\toprule
\textbf{Field} & $\SU{3}_{\mathrm{gauge}}$ & $\SU{6}_L$ & $\SU{6}_R$
  & $\UU{1}_B$ & $R_{\mathrm{sf}}$ & $R_{\mathrm{ferm}}$ \\
\midrule
$Q$ & $\fund$ & $\fund$ & $\mathbf{1}$ & $+1$ & $1/2$ & $-1/2$ \\
$\widetilde{Q}$ & $\antifund$ & $\mathbf{1}$ & $\antifund$ & $-1$
  & $1/2$ & $-1/2$ \\
$\lambda$ & $\adj$ & $\mathbf{1}$ & $\mathbf{1}$ & $0$ & $+1$ & $+1$ \\
\bottomrule
\end{tabular}
\end{center}

\textbf{Magnetic theory} ($\SU{3}_{\mathrm{mag}}$ gauge, $\Nf = 6$):

\begin{center}
\begin{tabular}{@{}l c c c c c c@{}}
\toprule
\textbf{Field} & $\SU{3}_{\mathrm{mag}}$ & $\SU{6}_L$ & $\SU{6}_R$
  & $\UU{1}_B$ & $R_{\mathrm{sf}}$ & $R_{\mathrm{ferm}}$ \\
\midrule
$q$ & $\fund$ & $\antifund$ & $\mathbf{1}$ & $+1$
  & $1/2$ & $-1/2$ \\
$\widetilde{q}$ & $\antifund$ & $\mathbf{1}$ & $\fund$ & $-1$
  & $1/2$ & $-1/2$ \\
$M$ & $\mathbf{1}$ & $\fund$ & $\antifund$ & $0$
  & $1$ & $0$ \\
$\lambda_{\mathrm{mag}}$ & $\adj$ & $\mathbf{1}$ & $\mathbf{1}$
  & $0$ & $+1$ & $+1$ \\
\bottomrule
\end{tabular}
\end{center}

Note: $B[q] = \Nc/(\Nf - \Nc) = 3/3 = 1$, which coincidentally equals
$B[Q]$ for this particular case.

\bigskip

\textbf{(b) Anomaly coefficient computations.}

\medskip

\textbf{(i) $\SU{6}_L^3$.}

\textit{Electric:} Only $Q$ transforms under $\SU{6}_L$
(as $\fund$), with $3$ gauge colours:
\[
  \mathcal{A}^{\mathrm{el}}_1 = 3 \times A(\fund)
  = 3 \times 1 = \boxed{3}\,.
\]

\textit{Magnetic:}  Two fields contribute:
\begin{itemize}
  \item $q$: $3$ mag.\ colours, $\antifund$ under $\SU{6}_L$:
    $3 \times A(\antifund) = 3 \times (-1) = -3$.
  \item $M$: gauge singlet, $\fund$ under $\SU{6}_L$ with $6$ copies
    (from the $\SU{6}_R$ antifund index):
    $6 \times A(\fund) = 6 \times 1 = 6$.
\end{itemize}
\[
  \mathcal{A}^{\mathrm{mag}}_1 = -3 + 6 = \boxed{3}\,.
  \quad\checkmark
\]

\medskip

\textbf{(ii) $\SU{6}_R^3$.}

By left--right symmetry ($Q \leftrightarrow \widetilde{Q}$,
$q \leftrightarrow \widetilde{q}$, $M^i{}_j \to M^j{}_i$):

\textit{Electric:} $\widetilde{Q}$ in $\antifund$ of $\SU{6}_R$,
$3$ colours: $3 \times A(\antifund) = -3$.

\textit{Magnetic:} $\widetilde{q}$ in $\fund$ of $\SU{6}_R$, $3$ mag
colours: $3 \times A(\fund) = 3$. $M$ in $\antifund$ of $\SU{6}_R$
with $6$ copies: $6 \times A(\antifund) = -6$. Total: $3 - 6 = -3$.

\[
  \mathcal{A}^{\mathrm{el}}_2 = -3 = \mathcal{A}^{\mathrm{mag}}_2\,.
  \quad\checkmark
\]

\medskip

\textbf{(iii) $\SU{6}_L^2\,\UU{1}_B$.}

Formula: $\mathcal{A} = \sum_f (\text{gauge mult}) \times T(R_L) \times
B_f$.

\textit{Electric:} $Q$: $3$ colours, $\fund$ of $\SU{6}_L$, $B = +1$:
\[
  \mathcal{A}^{\mathrm{el}}_3 = 3 \times \tfrac{1}{2} \times 1
  = \boxed{\tfrac{3}{2}}\,.
\]

\textit{Magnetic:}
\begin{itemize}
  \item $q$: $3$ colours, $\antifund$ of $\SU{6}_L$, $B = 1$:
    $3 \times \frac{1}{2} \times 1 = \frac{3}{2}$.
  \item $M$: $B = 0$, no contribution.
\end{itemize}
\[
  \mathcal{A}^{\mathrm{mag}}_3 = \boxed{\tfrac{3}{2}}\,.
  \quad\checkmark
\]

\medskip

\textbf{(iv) $\SU{6}_L^2\,\UU{1}_R$.}

Formula: $\mathcal{A} = \sum_f (\text{gauge mult}) \times T(R_L)
\times R_{\mathrm{ferm},f}$.

\textit{Electric:} $Q$: $3$ colours, $\fund$ of $\SU{6}_L$,
$R_{\mathrm{ferm}} = -1/2$:
\[
  \mathcal{A}^{\mathrm{el}}_4 = 3 \times \tfrac{1}{2} \times (-\tfrac{1}{2})
  = \boxed{-\tfrac{3}{4}}\,.
\]

\textit{Magnetic:}
\begin{itemize}
  \item $q$: $3$ mag colours, $\antifund$ of $\SU{6}_L$,
    $R_{\mathrm{ferm}} = -1/2$:
    $3 \times \frac{1}{2} \times (-\frac{1}{2}) = -\frac{3}{4}$.
  \item $M$: $6$ copies, $\fund$ of $\SU{6}_L$,
    $R_{\mathrm{ferm}} = 0$:
    $6 \times \frac{1}{2} \times 0 = 0$.
\end{itemize}
\[
  \mathcal{A}^{\mathrm{mag}}_4 = -\tfrac{3}{4} + 0
  = \boxed{-\tfrac{3}{4}}\,.
  \quad\checkmark
\]

\medskip

\textbf{(v) $\UU{1}_B^2\,\UU{1}_R$.}

Formula: $\mathcal{A} = \sum_f n_f \times B_f^2 \times R_{\mathrm{ferm},f}$,
where $n_f$ is the total multiplicity (gauge $\times$ flavour).

\textit{Electric:}
\begin{itemize}
  \item $Q$: $3 \times 6 = 18$ fermions, $B = 1$, $R_{\mathrm{ferm}} = -1/2$:
    $18 \times 1 \times (-1/2) = -9$.
  \item $\widetilde{Q}$: $18$ fermions, $B = -1$,
    $R_{\mathrm{ferm}} = -1/2$:
    $18 \times 1 \times (-1/2) = -9$.
  \item Gaugino: $B = 0$, no contribution.
\end{itemize}
\[
  \mathcal{A}^{\mathrm{el}}_5 = -9 + (-9) = \boxed{-18}\,.
\]

\textit{Magnetic:}
\begin{itemize}
  \item $q$: $3 \times 6 = 18$ fermions, $B = 1$, $R_{\mathrm{ferm}} = -1/2$:
    $18 \times 1 \times (-1/2) = -9$.
  \item $\widetilde{q}$: $18$ fermions, $B = -1$,
    $R_{\mathrm{ferm}} = -1/2$:
    $18 \times 1 \times (-1/2) = -9$.
  \item $M$: $B = 0$, no contribution.
  \item Gaugino: $B = 0$, no contribution.
\end{itemize}
\[
  \mathcal{A}^{\mathrm{mag}}_5 = -9 + (-9) = \boxed{-18}\,.
  \quad\checkmark
\]

\medskip

\textbf{(vi) $\Tr[R_{\mathrm{ferm}}]$ (gravitational).}

\textit{Electric:}
\begin{itemize}
  \item $Q$: $18$ fermions, $R_{\mathrm{ferm}} = -1/2$: $18 \times
    (-1/2) = -9$.
  \item $\widetilde{Q}$: $18$ fermions, $R_{\mathrm{ferm}} = -1/2$:
    $-9$.
  \item Gaugino: $3^2 - 1 = 8$ fermions, $R_{\mathrm{ferm}} = +1$:
    $+8$.
\end{itemize}
\[
  \mathcal{A}^{\mathrm{el}}_6 = -9 - 9 + 8 = \boxed{-10}\,.
\]

\textit{Magnetic:}
\begin{itemize}
  \item $q$: $3 \times 6 = 18$ fermions, $R_{\mathrm{ferm}} = -1/2$:
    $-9$.
  \item $\widetilde{q}$: $18$ fermions, $R_{\mathrm{ferm}} = -1/2$:
    $-9$.
  \item $M$: $6^2 = 36$ fermions, $R_{\mathrm{ferm}} = 0$:
    $0$.
  \item Gaugino: $3^2 - 1 = 8$ fermions, $R_{\mathrm{ferm}} = +1$:
    $+8$.
\end{itemize}
\[
  \mathcal{A}^{\mathrm{mag}}_6 = -9 - 9 + 0 + 8 = \boxed{-10}\,.
  \quad\checkmark
\]

\bigskip

\textbf{(c) Summary table.}

\begin{center}
\renewcommand{\arraystretch}{1.3}
\begin{tabular}{@{}l c c c@{}}
\toprule
\textbf{Coefficient} & \textbf{Electric} & \textbf{Magnetic} & \textbf{Match?} \\
\midrule
$\SU{6}_L^3$ & $3$ & $3$ & $\checkmark$ \\
$\SU{6}_R^3$ & $-3$ & $-3$ & $\checkmark$ \\
$\SU{6}_L^2\,\UU{1}_B$ & $3/2$ & $3/2$ & $\checkmark$ \\
$\SU{6}_L^2\,\UU{1}_R$ & $-3/4$ & $-3/4$ & $\checkmark$ \\
$\UU{1}_B^2\,\UU{1}_R$ & $-18$ & $-18$ & $\checkmark$ \\
$\Tr[R]$ (gravitational) & $-10$ & $-10$ & $\checkmark$ \\
\bottomrule
\end{tabular}
\end{center}

All six independent anomaly coefficients match.  These values are
consistent with the output of
\texttt{scripts/verify\_anomaly\_matching.py} for $\Nc = 3$,
$\Nf = 6$.
\end{proof}

\newpage


%% ========================================================================
%%  PROBLEM 2: Mass Deformation from N_f = 6 to N_f = 5
%% ========================================================================
\section*{Problem 2: Mass Deformation from $\Nf = 6$ to $\Nf = 5$
  \hfill $\star\star\star$ \normalsize(20 min)}
\addcontentsline{toc}{section}{Problem 2: Mass deformation
  $\Nf = 6 \to 5$}
\label{prob:mass-deformation}

\noindent\textbf{Background.}
A powerful consistency check of Seiberg duality is that giving a mass
to one flavour on both sides of the duality produces the correct dual
pair with $\Nf - 1$ flavours.  This was presented in
Lecture~3, \S5, for general $(\Nc, \Nf)$.

\medskip

\noindent\textbf{Questions.}

Start with $\SU{3}$ SQCD with $\Nf = 6$ flavours.  Add a mass term for
the 6th flavour: $\Delta\Wsup = m\, Q_6\, \widetilde{Q}^6$.

\begin{enumerate}[label=(\alph*)]
  \item What is the resulting \textbf{electric} theory after integrating
    out the massive quarks?

  \item What is the \textbf{magnetic} dual of this deformed theory?
    (Use the general Seiberg duality rule.)

  \item On the magnetic side of the \emph{original} $(\Nc = 3, \Nf = 6)$
    theory, the mass deformation maps to
    $\Delta\Wsup = m\, M^6{}_6$ (via $M^6{}_6 = Q_6\, \widetilde{Q}^6$).
    Show that the $F$-term equation for $M^6{}_6$ forces
    $\langle q_6\, \widetilde{q}^6\rangle \neq 0$, and explain how this
    \emph{Higgses} the magnetic gauge group from $\SU{3}$ to $\SU{2}$.

  \item As a check, verify the anomaly coefficient $\SU{5}_L^3$ for the
    deformed dual pair $(\SU{3}, \Nf = 5) \leftrightarrow (\SU{2},
    \Nf = 5)$.
\end{enumerate}

\medskip

\noindent\textit{Hint for (c):}  The total magnetic superpotential is
$\Wsup = (1/\mu)\, M^i{}_j\, q_i\, \widetilde{q}^j + m\, M^6{}_6$.
Setting $\partial\Wsup/\partial M^6{}_6 = 0$ gives
$q_6\,\widetilde{q}^6 = -\mu\, m$.

\bigskip
\hrule
\bigskip

\begin{proof}[\textbf{Solution to Problem 2}]

\textbf{(a)} The mass term $m\, Q_6\,\widetilde{Q}^6$ makes the 6th
flavour heavy.  Below the scale $m$, integrating out $Q_6$ and
$\widetilde{Q}^6$ leaves:
\[
  \boxed{\text{Electric:}\quad \SU{3} \text{ with } \Nf = 5
  \text{ flavours, } \Wsup = 0\,.}
\]

\textbf{(b)} Applying the Seiberg duality rule to
$\SU{3}$ with $\Nf = 5$:
\[
  \text{Magnetic gauge group:}\quad \SU{\Nf - \Nc} = \SU{5 - 3} = \SU{2}\,,
\]
with $5$ magnetic quark pairs $(q_i, \widetilde{q}^i)$, a $5 \times 5$
meson matrix $M'$, and $\Wsup_{\mathrm{mag}} = (1/\mu')\, M' q \widetilde{q}$.

\[
  \boxed{\text{Magnetic:}\quad \SU{2} \text{ with } \Nf = 5
  \text{ flavours } + M' + \Wsup_{\mathrm{mag}}\,.}
\]

\textbf{(c)} The total superpotential of the magnetic $\SU{3}$ theory
with the mass deformation is:
\[
  \Wsup = \frac{1}{\mu}\, M^i{}_j\, q_i\, \widetilde{q}^j + m\, M^6{}_6\,.
\]
The $F$-term equation for $M^6{}_6$ is:
\[
  \frac{\partial\Wsup}{\partial M^6{}_6} = 0
  \qquad\Longrightarrow\qquad
  \frac{1}{\mu}\, q_6\, \widetilde{q}^6 + m = 0
  \qquad\Longrightarrow\qquad
  \langle q_6\, \widetilde{q}^6\rangle = -\mu\, m \neq 0\,.
\]
The nonzero VEV of $q_6$ (in the fundamental of $\SU{3}_{\mathrm{mag}}$)
breaks the magnetic gauge group by the Higgs mechanism:
\[
  \SU{3}_{\mathrm{mag}}
  \;\xrightarrow{\;\langle q_6 \rangle \neq 0\;}
  \;\SU{2}_{\mathrm{mag}}\,.
\]
After integrating out the massive fields ($q_6$, $\widetilde{q}^6$,
$M^6{}_j$, $M^i{}_6$), the surviving magnetic theory is precisely
$\SU{2}$ with $5$ flavours plus the $5 \times 5$ meson $M'$, matching
the prediction of part~(b).\quad$\checkmark$

\textbf{(d)} Check $\SU{5}_L^3$ for $(\SU{3}, \Nf = 5) \leftrightarrow
(\SU{2}, \Nf = 5)$.

\textit{Electric:} $Q$ in $\fund$ of $\SU{5}_L$, $3$ colours:
$\mathcal{A}^{\mathrm{el}} = 3 \times 1 = 3$.

\textit{Magnetic:} Magnetic gauge group is $\SU{2}$.
\begin{itemize}
  \item $q$: $2$ mag colours, $\antifund$ of $\SU{5}_L$:
    $2 \times (-1) = -2$.
  \item $M'$: $5$ copies, $\fund$ of $\SU{5}_L$:
    $5 \times 1 = 5$.
\end{itemize}
$\mathcal{A}^{\mathrm{mag}} = -2 + 5 = 3$.

\[
  \boxed{\mathcal{A}^{\mathrm{el}} = 3 = \mathcal{A}^{\mathrm{mag}}}\,.
  \quad\checkmark
\]
\end{proof}

\newpage


%% ========================================================================
%%  PROBLEM 3: R-charges and the Superpotential
%% ========================================================================
\section*{Problem 3: $R$-charges and the Superpotential
  \hfill $\star\star$ \normalsize(15 min)}
\addcontentsline{toc}{section}{Problem 3: $R$-charges and the
  superpotential}
\label{prob:r-charges}

\noindent\textbf{Background.}
The non-anomalous $R$-symmetry plays a central role in Seiberg duality.
At the superconformal fixed point, the $R$-charge determines the
scaling dimension of chiral primary operators via
$\Delta = \tfrac{3}{2}\,|R|$.

\medskip

\noindent\textbf{Questions.}
\begin{enumerate}[label=(\alph*)]
  \item \textbf{Derive $\Rcharge{Q}$.}  For $\SU{\Nc}$ SQCD with $\Nf$
    massless flavours, the anomaly-free $R$-symmetry satisfies
    $\Tr[\SU{\Nc}^2\,\UU{1}_R] = 0$.  Using the fermion $R$-charges
    and the fact that the gaugino has $R_{\mathrm{ferm}} = +1$, derive:
    \[
      \Rcharge{Q} = \frac{\Nf - \Nc}{\Nf}\,.
    \]

  \item \textbf{Magnetic $R$-charges.}  Using the Seiberg duality
    map and the requirement $\Rcharge{\Wsup_{\mathrm{mag}}} = 2$,
    determine $\Rcharge{q}$, $\Rcharge{\widetilde{q}}$, and
    $\Rcharge{M}$ as functions of $\Nc$ and $\Nf$.

  \item \textbf{Verify $\Rcharge{\Wsup_{\mathrm{mag}}} = 2$.}  Show that
    \[
      \Rcharge{M} + \Rcharge{q} + \Rcharge{\widetilde{q}} = 2\,.
    \]

  \item \textbf{Conformal dimension and the conformal window.}
    At the IR fixed point, the conformal dimension of a chiral primary
    scalar operator $\mathcal{O}$ is $\Delta(\mathcal{O}) =
    \tfrac{3}{2}\,R[\mathcal{O}]$.  Compute the dimension of the
    meson $M = Q\widetilde{Q}$ in the electric theory.  For which value
    of $\Nf$ does $\Delta[M] = 1$ (free field)?  What is the physical
    significance of this value?
\end{enumerate}

\medskip

\noindent\textit{Hint for (a):} The anomaly coefficient
$\Tr[\SU{\Nc}^2\,\UU{1}_R]$ receives contributions from all fermions
charged under the gauge group: $Q$-fermions, $\widetilde{Q}$-fermions,
and the gaugino.  Each contributes $T(R) \times R_{\mathrm{ferm}}$.

\bigskip
\hrule
\bigskip

\begin{proof}[\textbf{Solution to Problem 3}]

\textbf{(a)} The anomaly $\Tr[\SU{\Nc}^2\,\UU{1}_R] = 0$ requires
the sum over all gauge-charged fermions:
\begin{align}
  0 &= \Nf \cdot T(\fund) \cdot R_{\mathrm{ferm}}[Q]
  + \Nf \cdot T(\antifund) \cdot R_{\mathrm{ferm}}[\widetilde{Q}]
  + T(\adj) \cdot R_{\mathrm{ferm}}[\lambda] \nonumber \\
  &= \Nf \cdot \tfrac{1}{2} \cdot (R[Q] - 1)
  + \Nf \cdot \tfrac{1}{2} \cdot (R[Q] - 1)
  + \Nc \cdot 1 \nonumber \\
  &= \Nf\,(R[Q] - 1) + \Nc\,.
\end{align}
Solving for $R[Q]$:
\[
  \Nf\, R[Q] = \Nf - \Nc
  \qquad\Longrightarrow\qquad
  \boxed{R[Q] = \frac{\Nf - \Nc}{\Nf}}\,.
\]

\textbf{(b)} The meson field $M$ maps to the composite $Q\widetilde{Q}$,
so its $R$-charge is:
\[
  R[M] = R[Q] + R[\widetilde{Q}] = 2\,\frac{\Nf - \Nc}{\Nf}\,.
\]
From the requirement $R[W_{\mathrm{mag}}] = R[M] + R[q] + R[\widetilde{q}] = 2$
and the symmetry $R[q] = R[\widetilde{q}]$:
\[
  R[q] = R[\widetilde{q}] = \frac{2 - R[M]}{2}
  = \frac{2 - 2(\Nf - \Nc)/\Nf}{2}
  = \frac{\Nc}{\Nf}\,.
\]
\[
  \boxed{R[q] = R[\widetilde{q}] = \frac{\Nc}{\Nf}\,,
  \qquad R[M] = \frac{2(\Nf - \Nc)}{\Nf}\,.}
\]

\textbf{(c)} Verification:
\[
  R[M] + R[q] + R[\widetilde{q}]
  = \frac{2(\Nf - \Nc)}{\Nf} + \frac{\Nc}{\Nf} + \frac{\Nc}{\Nf}
  = \frac{2\Nf - 2\Nc + 2\Nc}{\Nf}
  = \frac{2\Nf}{\Nf} = 2\,.
  \quad\checkmark
\]

\textbf{(d)} The conformal dimension of the meson is:
\begin{align}
  \Delta[M] &= \tfrac{3}{2}\,R[M]
  = \tfrac{3}{2} \times \frac{2(\Nf - \Nc)}{\Nf}
  = \frac{3(\Nf - \Nc)}{\Nf}\,.
\end{align}
Setting $\Delta[M] = 1$ (the dimension of a free scalar field):
\begin{align}
  \frac{3(\Nf - \Nc)}{\Nf} &= 1\,, \\
  3\Nf - 3\Nc &= \Nf\,, \\
  2\Nf &= 3\Nc\,, \\
  \Nf &= \frac{3\Nc}{2}\,.
\end{align}

\[
  \boxed{\Delta[M] = 1 \quad\text{at}\quad \Nf = \frac{3\Nc}{2}\,.}
\]

\textbf{Physical significance:} $\Nf = 3\Nc/2$ is the \emph{lower
boundary of the SUSY conformal window}.  At this value, the meson
becomes a free field and decouples from the interacting sector.  Below
$\Nf = 3\Nc/2$, the conformal dimension would fall below the unitarity
bound $\Delta \geq 1$, signalling that the conformal description breaks
down.  The magnetic theory enters the free magnetic phase: it is
IR-free, and the correct description is in terms of the weakly-coupled
magnetic quarks and mesons.

This connects directly to Exercise Set~2, Problem~3(c), where the
conformal window was determined from the unitarity bound.
\end{proof}

\newpage


%% ========================================================================
%%  PROBLEM 4: s-Confinement at N_f = N_c + 1
%% ========================================================================
\section*{Problem 4: $s$-Confinement at $\Nf = \Nc + 1$
  \hfill $\star\star\star$ \normalsize(25 min)}
\addcontentsline{toc}{section}{Problem 4: $s$-Confinement at
  $\Nf = \Nc + 1$}
\label{prob:s-confinement}

\noindent\textbf{Background.}
The case $\Nf = \Nc + 1$ is a limiting case of Seiberg duality where the
magnetic gauge group becomes trivial.  The theory \emph{s-confines}:
it confines without chiral symmetry breaking, and the low-energy
description consists entirely of gauge-singlet composite fields.  This
was first analysed by Seiberg~\cite{Seiberg1995} and systematically
by Intriligator and Seiberg~\cite{IS1995}.

\medskip

\noindent\textbf{Questions.}
\begin{enumerate}[label=(\alph*)]
  \item In the Seiberg dual of $\SU{\Nc}$ SQCD with $\Nf = \Nc + 1$
    flavours, what is the magnetic gauge group?

  \item Since the magnetic gauge group is trivial, the magnetic ``quarks''
    $q_i$ and $\widetilde{q}^j$ are gauge-singlet fields.  These are
    identified with the electric baryons:
    \[
      B_i \;\longleftrightarrow\; q_i\,,
      \qquad
      \widetilde{B}^j \;\longleftrightarrow\; \widetilde{q}^j\,.
    \]
    Together with the mesons $M^i{}_j$, these form the complete
    set of light degrees of freedom.  Write the most general
    superpotential consistent with the global symmetries
    $\SU{\Nf}_L \times \SU{\Nf}_R \times \UU{1}_B \times \UU{1}_R$.

  \item Show that the exact confining superpotential takes the form:
    \begin{equation}\label{eq:W-sconf}
      \Wsup = \frac{1}{\Lambda^{2\Nc - 1}}
      \left(B_i\, M^i{}_j\, \widetilde{B}^j - \det M\right).
    \end{equation}
    Verify the \textbf{dimensional consistency}: show that each term
    has dimension $[\mathrm{mass}]^3$.

  \item For $\Nc = 3$, $\Nf = 4$, verify the anomaly coefficient
    $\SU{4}_L^3$ for the confined description (mesons $M$, baryons
    $B_i$, $\widetilde{B}^j$) against the fundamental description
    ($Q$, $\widetilde{Q}$, $\lambda$).
\end{enumerate}

\medskip

\noindent\textit{Hints:}

(a) Apply $\SU{\Nf - \Nc} = \SU{(\Nc+1) - \Nc} = \SU{1}$.

(c) The mesons $M$ have dimension $[\mathrm{mass}]$ (see CONVENTIONS.md).
The baryons $B$ are products of $\Nc$ fundamental fields, each of
dimension $[\mathrm{mass}]$, so $[B] = [\mathrm{mass}]^{\Nc}$.
The dynamical scale $\Lambda$ has $[\Lambda] = [\mathrm{mass}]$.

(d) The baryon $B_i$ transforms as $\antifund$ of $\SU{\Nf}_L$
(antisymmetric product of $\Nc$ fundamentals with $\Nf = \Nc + 1$
gives the antifundamental of $\SU{\Nc+1}_L$).

\bigskip
\hrule
\bigskip

\begin{proof}[\textbf{Solution to Problem 4}]

\textbf{(a)} The magnetic gauge group is:
\[
  \SU{\Nf - \Nc} = \SU{(\Nc + 1) - \Nc} = \SU{1}\,.
\]
But $\SU{1}$ is the trivial group (containing only the identity).
Therefore:
\[
  \boxed{\text{There is \textbf{no} magnetic gauge group.}}
\]
The magnetic theory has no gauge dynamics: all ``magnetic'' fields are
gauge singlets.

\textbf{(b)} The confined degrees of freedom and their quantum numbers
under $\SU{\Nf}_L \times \SU{\Nf}_R \times \UU{1}_B \times \UU{1}_R$:

\renewcommand{\arraystretch}{1.3}
\begin{center}
\begin{tabular}{@{}l c c c c@{}}
\toprule
\textbf{Field} & $\SU{\Nf}_L$ & $\SU{\Nf}_R$ & $\UU{1}_B$ & $R_{\mathrm{sf}}$ \\
\midrule
$M^i{}_j$ & $\fund$ & $\antifund$ & $0$
  & $2(\Nf - \Nc)/\Nf = 2/(\Nc + 1)$ \\
$B_i$ & $\antifund$ & $\mathbf{1}$ & $+\Nc$
  & $\Nc(\Nf - \Nc)/\Nf = \Nc/(\Nc + 1)$ \\
$\widetilde{B}^j$ & $\mathbf{1}$ & $\fund$ & $-\Nc$
  & $\Nc/(\Nc + 1)$ \\
\bottomrule
\end{tabular}
\end{center}

(Here $B_i = \epsilon^{a_1\cdots a_{\Nc}}\, Q^{i_1}_{a_1} \cdots
Q^{i_{\Nc}}_{a_{\Nc}}$ transforms as the $\Nc$-th antisymmetric product
of $\fund$ under $\SU{\Nf}_L$; for $\Nf = \Nc + 1$, this is the
$\antifund$ of $\SU{\Nc+1}_L$.)

The most general superpotential consistent with these symmetries is:
\begin{equation}
  \Wsup = \frac{\alpha}{\Lambda^{2\Nc - 1}}\, B_i\, M^i{}_j\,
  \widetilde{B}^j
  + \frac{\beta}{\Lambda^{2\Nc - 1}}\, \det M\,,
\end{equation}
where $\alpha$ and $\beta$ are dimensionless constants.  These are the
only two independent invariants that can be formed:
\begin{itemize}
  \item $B_i\, M^i{}_j\, \widetilde{B}^j$ contracts all flavour
    indices and has the correct $R$-charge, baryon charge, and flavour
    transformation properties.
  \item $\det M$ is an $\SU{\Nf}_L \times \SU{\Nf}_R$ invariant with
    $B = 0$ and $R = \Nf \times 2/(\Nc+1) = 2$ (using
    $\Nf = \Nc + 1$).
\end{itemize}

\textbf{(c)} The exact result from~\cite{Seiberg1995,IS1995} is
$\alpha = 1$, $\beta = -1$:
\[
  \boxed{
    \Wsup = \frac{1}{\Lambda^{2\Nc - 1}}
    \left(B_i\, M^i{}_j\, \widetilde{B}^j - \det M\right).
  }
\]
(The relative sign and normalisation are fixed by requiring smooth
limits as various masses are taken to infinity, matching the ADS
superpotential for $\Nf < \Nc$, and by instanton calculations.)

\textbf{Dimensional check.}  We need
$[\Wsup] = [\mathrm{mass}]^3$:
\begin{itemize}
  \item $[M] = [\mathrm{mass}]$ (elementary meson field).
  \item $[B] = [\mathrm{mass}]^{\Nc}$ (product of $\Nc$ fundamentals,
    each of dimension $[\mathrm{mass}]$).
  \item $[\Lambda] = [\mathrm{mass}]$.
\end{itemize}

For the first term:
\[
  \frac{[B \cdot M \cdot \widetilde{B}]}{[\Lambda^{2\Nc-1}]}
  = \frac{[\mathrm{mass}]^{\Nc} \cdot [\mathrm{mass}]
  \cdot [\mathrm{mass}]^{\Nc}}{[\mathrm{mass}]^{2\Nc - 1}}
  = \frac{[\mathrm{mass}]^{2\Nc + 1}}{[\mathrm{mass}]^{2\Nc - 1}}
  = [\mathrm{mass}]^2\,.
\]

Wait --- this gives $[\mathrm{mass}]^2$, not $[\mathrm{mass}]^3$.  We
need to be more careful about the dimension of $B$.

\textbf{Corrected dimensional analysis.}  In SQCD, the elementary chiral
superfield $Q$ has classical dimension $[Q] = [\mathrm{mass}]$ (the scalar
component of a chiral superfield has mass dimension $1$ in $d=4$).
The baryon operator is:
\[
  B = \epsilon^{a_1\cdots a_{\Nc}}\, Q_{a_1}^{i_1} \cdots
  Q_{a_{\Nc}}^{i_{\Nc}}\,,
  \qquad [B] = [\mathrm{mass}]^{\Nc}\,.
\]
The meson $M^i{}_j = Q^i \widetilde{Q}_j$ has $[M] = [\mathrm{mass}]^2$
in the \emph{electric} theory (it is a composite of two dimension-1
fields).

In the \emph{magnetic} (confined) description, $M$ is an elementary
field with $[M] = [\mathrm{mass}]^2$ inherited from the duality map.
(Alternatively, at the conformal fixed point, the dimension is
$\Delta[M] = (3/2)R[M]$; but for the superpotential counting, we use
the UV dimension.)

Now:
\[
  \frac{[B \cdot M \cdot \widetilde{B}]}{[\Lambda^{2\Nc-1}]}
  = \frac{[\mathrm{mass}]^{\Nc} \cdot [\mathrm{mass}]^2
  \cdot [\mathrm{mass}]^{\Nc}}{[\mathrm{mass}]^{2\Nc - 1}}
  = \frac{[\mathrm{mass}]^{2\Nc + 2}}{[\mathrm{mass}]^{2\Nc - 1}}
  = [\mathrm{mass}]^3\,.
  \quad\checkmark
\]

For $\det M$ with $\Nf = \Nc + 1$:
\[
  \frac{[\det M]}{[\Lambda^{2\Nc-1}]}
  = \frac{[\mathrm{mass}]^{2\Nf}}{[\mathrm{mass}]^{2\Nc-1}}
  = \frac{[\mathrm{mass}]^{2(\Nc+1)}}{[\mathrm{mass}]^{2\Nc-1}}
  = [\mathrm{mass}]^{2\Nc + 2 - 2\Nc + 1}
  = [\mathrm{mass}]^3\,.
  \quad\checkmark
\]

Both terms have the correct dimension.

\bigskip

\textbf{(d)} Anomaly check for $\SU{4}_L^3$ with $\Nc = 3$,
$\Nf = 4$.

\textit{Electric (fundamental) theory:}
\[
  Q \text{ in } \fund \text{ of } \SU{4}_L\,,\quad 3 \text{ colours:}\quad
  \mathcal{A}^{\mathrm{el}} = 3 \times A(\fund) = 3 \times 1 = 3\,.
\]

\textit{Confined theory:}
\begin{itemize}
  \item $M^i{}_j$: transforms as $\fund$ under $\SU{4}_L$, with $4$
    copies (from $\SU{4}_R$ antifund index):
    $4 \times A(\fund) = 4 \times 1 = 4$.
  \item $B_i$: transforms as $\antifund$ of $\SU{4}_L$ (the $3$-index
    antisymmetric product of $\fund$ of $\SU{4}$ is $\antifund$ of
    $\SU{4}$, since $\binom{4}{3} = 4$ and it transforms as the dual):
    $1 \times A(\antifund) = 1 \times (-1) = -1$.
  \item $\widetilde{B}^j$: singlet under $\SU{4}_L$, no contribution.
\end{itemize}
\[
  \mathcal{A}^{\mathrm{conf}} = 4 + (-1) + 0 = 3\,.
\]

\[
  \boxed{\mathcal{A}^{\mathrm{el}} = 3
  = \mathcal{A}^{\mathrm{conf}}}\,.
  \quad\checkmark
\]

\begin{remark}[The s-confinement programme]
  The result $\SU{\Nc}$ with $\Nf = \Nc + 1$ is the simplest case of
  \emph{s-confinement} (confinement without chiral symmetry breaking).
  The systematic classification of all s-confining theories was carried
  out by Csaki, Schmaltz, and Skiba~\cite{CSS1997}.  The existence of an
  exact confining superpotential~\eqref{eq:W-sconf} is one of the most
  striking predictions of supersymmetric dynamics.
\end{remark}

\end{proof}

\bigskip

\begin{center}
\rule{0.6\textwidth}{0.4pt}

\textit{End of Exercise Set 3.  These exercises provide hands-on
practice with the core machinery of Seiberg duality: anomaly matching,
mass deformations, $R$-charge analysis, and s-confinement.
The tools developed here will be essential when we apply duality
to the Standard Model in Lectures~5--7.}
\end{center}


%% ========================================================================
%%  BIBLIOGRAPHY
%% ========================================================================
\begin{thebibliography}{99}

\bibitem{Seiberg1995}
  N.~Seiberg,
  ``Electric--magnetic duality in supersymmetric non-Abelian gauge
  theories,''
  \textit{Nucl.~Phys.} \textbf{B435} (1995) 129
  [\texttt{hep-th/9411149}].

\bibitem{IS1995}
  K.~Intriligator and N.~Seiberg,
  ``Lectures on supersymmetric gauge theories and electric-magnetic
  duality,''
  \textit{Nucl.~Phys.~Proc.~Suppl.} \textbf{45BC} (1996) 1
  [\texttt{hep-th/9509066}].

\bibitem{CSS1997}
  C.~Cs\'aki, M.~Schmaltz, and W.~Skiba,
  ``A systematic approach to confinement in $\mathcal{N}=1$
  supersymmetric gauge theories,''
  \textit{Phys.~Rev.~Lett.} \textbf{78} (1997) 799
  [\texttt{hep-th/9610139}].

\bibitem{Sannino2024}
  F.~Sannino \textit{et al.},
  ``The Dualised Standard Model,''
  [\texttt{arXiv:2407.17281}].

\end{thebibliography}

\end{document}
